\begin{abstract}
Optimizing through a fixed neural surrogate couples two computational systems:
automatic differentiation through a network and a constrained nonlinear
programming solver.  Moving the first to a GPU does not imply that the
Karush--Kuhn--Tucker (KKT) system should follow, nor that exact curvature is the
right model.  We present \method, an audited JAX--MadNLP interface that exposes
neural derivatives through DLPack device buffers while preserving
problem-supplied Jacobian and Lagrangian-Hessian sparsity.  We isolate four
choices---structure disclosure, curvature, transfer semantics, and KKT
placement---on controlled problems and three established neural-embedded NLP
families.  Correct sparsity reduces a published neural NMPC solve from
\PFRCaseOneDenseCPU\,s to \PFRCaseOneSparseCPU\,ms and reverses the CPU/GPU
ranking; at 7,500 neural equalities, however, GPU KKT is
\PFRCaseThreeSpeedup$\times$ faster than CPU KKT.  A prospectively evaluated
curvature route raises Darcy inversion success from \DarcyRouteBase/\DarcyRouteTotal{}
to \DarcyRouteTotal/\DarcyRouteTotal{} at \DarcyRouteOverhead\% additional warm
time.  At a fixed FNO parameter budget of 1.18M, changing depth and width moves the
exact-Hessian peak from \FNOShapeMinMemory{} to \FNOShapeMaxMemory{} GiB;
parameter count alone does not predict second-order AD memory.  Finally,
independent PDE replay shows that surrogate improvement need
not transfer: Gauss--Newton lowers FNO residual relative to L-BFGS on all
twelve physical instances, while the paired PDE change has a 95\% bootstrap
interval spanning $\ReplayPairedPDEDeltaLow$ to
$\ReplayPairedPDEDeltaHigh$ percentage points.  Descriptive start-level
rankings agree in only \ReplayAgreement/\ReplayGroups{} matched groups.  The
result is a reproducible regime map rather than a single GPU speedup: solver
convergence, systems speed, and application validity remain distinct
requirements.
\end{abstract}
